\begin{table}[htbp]
\centering
\footnotesize
\captionstyle{\centering}
\normalcaption
\caption{Cobertura da conversão por predominância territorial, 2001--2024}\label{tab:h1-cobertura}
\begin{tabular}{lrrrrr}
\toprule
Grupo & Municípios & Mistos & Pares válidos & \% da grade & Fixos \\
\midrule
Amazônia & 76 & 35 & 1.306 & 71,6 & 39 \\
Cerrado & 61 & 24 & 1.412 & 96,4 & 55 \\
Pantanal & 4 & 4 & 54 & 56,2 & 1 \\
\bottomrule
\end{tabular}
\par\smallskip\begin{minipage}{\textwidth}\footnotesize
\SingleSpacing
\noindent\textbf{Fonte:} elaboração própria com MapBiomas 10.1 e IBGE.\par\medskip
\noindent\textbf{Nota:} Os 141 municípios (grade teórica de 3.384 pares em 24 anos) agrupam-se pelo bioma predominante; ``Mistos'' indica municípios que abrangem mais de um bioma (63). O Cerrado alcança 96,4\% de cobertura livre de nuvens (55 dos 95 municípios contínuos em ``Fixos''), contra 71,6\% na Amazônia e 56,2\% no Pantanal; as regressões retêm 2.767 dos 2.772 pares válidos após excluir grupos unitários.
\end{minipage}
\end{table}
